\chapter{Фреймворк Spring}
\label{ch:07}

В предыдущей главе мы научились собирать проект системами Maven и Gradle и
разговаривать с базой данных через JDBC и Hibernate. Но стоит соединить эти
части в настоящее веб-приложение~--- с контроллерами, сервисами, формой входа
и REST API,~--- как выясняется, что «склеивать» компоненты вручную крайне
неудобно. Эту инфраструктурную работу берёт на себя \term{Spring}~--- де-факто
стандарт корпоративной разработки на Java. Прочитав главу, вы разберётесь в
инверсии управления и внедрении зависимостей, научитесь описывать бины,
соберёте приложение со слоями «контроллер~--- сервис~--- репозиторий»,
отдадите данные через REST API и HTML-страницы, закроете его аутентификацией
и доведёте до состояния, пригодного для эксплуатации.

\section{Инверсия управления и внедрение зависимостей}

Представьте класс \texttt{StudentService}, который сохраняет студентов в базу.
Ему нужен \texttt{StudentRepository}, репозиторию~--- источник данных, а
источнику~--- адрес базы, логин и пароль. В классическом подходе сервис
добывает всё это сам: в его конструкторе стоят вызовы \texttt{new}, а имя
конкретного класса репозитория вписано намертво. Такой код порождает четыре
беды сразу: \emph{жёсткую связность}~--- заменить реализацию нельзя, не
переписав конструктор; \emph{нулевую тестируемость}~--- репозиторий не
подменить заглушкой, и тест полезет в настоящую базу; \emph{дублирование}~---
источник данных создаётся заново в каждом сервисе; \emph{ручной жизненный
цикл}~--- соединения и пулы вы обязаны закрывать сами.

\term{Инверсия управления} (\eng{Inversion of Control, IoC})~--- принцип
проектирования, при котором создание объектов, их связывание и управление
жизненным циклом передаются внешнему контейнеру. Формулируют его
«голливудским принципом»: «не звоните нам~--- мы сами вам позвоним». Бытовая
аналогия~--- ресторан: вы не идёте на кухню за кастрюлей и не ищете, где растёт
картофель, а говорите, что хотите получить. \term{Внедрение зависимостей}
(\eng{Dependency Injection, DI})~--- техника реализации IoC: объект не создаёт
свои зависимости, а получает их извне готовыми. Способов три, и все они
собраны в листинге~\ref{lst:07-1}.

\begin{listing}[H]
\begin{minted}{java}
@Service                    // 1. Через конструктор - рекомендуется
public class StudentService {
    private final StudentRepository repository;

    // Spring сам создаст репозиторий и передаст его сюда
    public StudentService(StudentRepository repository) {
        this.repository = repository;
    }
}

@Service                    // 2. Через сеттер
public class SetterDemo {
    private StudentRepository repository;

    @Autowired
    public void setRepository(StudentRepository repo) {
        this.repository = repo;
    }
}

@Service                    // 3. Через поле - не рекомендуется
public class FieldDemo {
    @Autowired
    private StudentRepository repository;
}
\end{minted}
\caption{Три способа внедрения зависимости}
\label{lst:07-1}
\end{listing}

Предпочтительно \term{конструкторное внедрение}. Только оно позволяет объявить
поле \texttt{final}, то есть гарантировать, что зависимость задана один раз;
если подходящего бина нет, приложение не поднимется вовсе~--- ошибка вылезет
при старте, а не через неделю в работающем сервисе; наконец, такой класс
тестируется без всякого Spring. Сеттер оставляют для необязательных
зависимостей. Внедрение в поле выглядит короче всех и именно поэтому опасно:
зависимости не видны в сигнатуре класса, а подставить заглушку без рефлексии
невозможно.

\section{Spring Framework и Spring Boot}

\term{Spring Framework}~--- универсальный фреймворк с открытым исходным кодом
для платформы Java, созданный, чтобы упростить разработку корпоративных
приложений: он управляет объектами, транзакциями и конфигурацией, а в основе
лежат IoC и DI. Важная черта Spring~--- ненавязчивость: он работает с обычными
Java-объектами и не требует наследоваться от специальных классов.
\term{Spring Boot}~--- надстройка над ним, позволяющая писать самодостаточные
приложения без единой строчки XML-конфигурации. Держится она на трёх китах.
\term{Автоконфигурация}: Spring осматривает classpath и сам подбирает разумные
настройки~--- увидел драйвер H2, значит, настроил источник данных.
\term{Встроенный сервер} (Tomcat, Jetty или Undertow): приложение запускается
обычным методом \texttt{main}. \term{Starter-модули}: одна зависимость
подключает набор согласованных библиотек.

Всё приложение Spring Boot умещается в шесть строк
(листинг~\ref{lst:07-2})~--- дальше остаётся только добавлять компоненты.

\begin{listing}[H]
\begin{minted}{java}
@SpringBootApplication
public class StudentApplication {
    public static void main(String[] args) {
        SpringApplication.run(StudentApplication.class, args);
    }
}
\end{minted}
\caption{Минимальное приложение Spring Boot}
\label{lst:07-2}
\end{listing}

Аннотация \texttt{@SpringBootApplication}~--- это «всё в одном»: она объединяет
\texttt{@Configuration} (класс может объявлять бины),
\texttt{@EnableAutoConfiguration} (включает автоконфигурацию) и
\texttt{@ComponentScan} (велит искать компоненты в этом пакете и всех
вложенных). Отсюда правило: главный класс кладут в корневой пакет проекта,
иначе часть классов останется вне сканирования. Заготовку проекта создают в
\term{Spring Initializr}~--- веб-сервисе по адресу
\texttt{start.spring.io}: там выбирают систему сборки, язык, версии и
зависимости. Самые нужные стартеры перечислены в табл.~\ref{tab:07-1}; полное
имя каждого начинается с \texttt{spring-boot-}. Подключают их как обычные
зависимости Maven, причём без указания версий~--- версии задаёт родительский
модуль \texttt{spring-boot-starter-parent}, и потому библиотеки не
конфликтуют.

\begin{table}[htbp]
\caption{Основные starter-модули Spring Boot}
\label{tab:07-1}
\centering
\begin{tabular}{L{5.4cm}L{9.4cm}}
\toprule
Стартер & Что подключает \\
\midrule
\texttt{starter-web} & Spring MVC и встроенный Tomcat \\
\texttt{starter-data-jpa} & работу с базой через JPA и Hibernate \\
\texttt{starter-security} & аутентификацию и авторизацию \\
\texttt{starter-thymeleaf} & HTML-шаблоны Thymeleaf \\
\texttt{starter-validation} & проверку данных (Hibernate Validator) \\
\texttt{starter-aop} & аспектно-ориентированное программирование \\
\bottomrule
\end{tabular}
\end{table}

Как это выглядит в файле \texttt{pom.xml}, показывает
листинг~\ref{lst:07-3}; остальные стартеры подключаются так же, меняется
только \texttt{artifactId}.

\begin{listing}[H]
\begin{minted}{xml}
<parent>
    <groupId>org.springframework.boot</groupId>
    <artifactId>spring-boot-starter-parent</artifactId>
    <version>3.5.4</version>
</parent>

<dependencies>
    <dependency>
        <groupId>org.springframework.boot</groupId>
        <artifactId>spring-boot-starter-web</artifactId>
    </dependency>
    <dependency>
        <groupId>com.h2database</groupId>
        <artifactId>h2</artifactId>
        <scope>runtime</scope>
    </dependency>
</dependencies>
\end{minted}
\caption{Описание проекта в файле pom.xml}
\label{lst:07-3}
\end{listing}

Настройки приложения хранятся в файле \texttt{application.properties} из
каталога \texttt{src/main/resources}. Набор для учебного проекта приведён в
листинге~\ref{lst:07-4}: база H2 живёт в оперативной памяти и
исчезает после остановки приложения, а Hibernate сам создаёт под сущности
таблицы.

\begin{listing}[H]
\begin{minted}{text}
spring.datasource.url=jdbc:h2:mem:testdb
spring.datasource.username=sa
spring.datasource.password=
# none / validate / update / create / create-drop
spring.jpa.hibernate.ddl-auto=update
# веб-консоль H2 открывается по пути /h2-console
spring.h2.console.enabled=true
spring.thymeleaf.cache=false
\end{minted}
\caption{Файл application.properties}
\label{lst:07-4}
\end{listing}

\section{Контейнер IoC и бины}

\term{Контейнер IoC}~--- ядро Spring: он создаёт объекты по метаданным
конфигурации, внедряет зависимости и управляет их жизненным циклом. Хранит
контейнер не сами объекты, а их описания~--- структуры
\texttt{BeanDefinition} с классом, зависимостями, областью видимости и
методами инициализации. Базовая форма контейнера~---
интерфейс \texttt{BeanFactory}, но на практике берут его наследника
\texttt{ApplicationContext}: тот добавляет интернационализацию, события, работу
с ресурсами и интеграцию со Spring AOP. В приложении Spring Boot контекст
создаётся вызовом \texttt{SpringApplication.run}.

\term{Бин} (\eng{bean})~--- объект, который создаёт, настраивает и хранит
контейнер. Живёт он по одному сценарию: контейнер читает конфигурацию и строит
\texttt{BeanDefinition}; создаёт экземпляр; внедряет в него зависимости;
вызывает метод инициализации (\texttt{@PostConstruct} или
\texttt{init-method}); выдаёт бин по запросу; при закрытии контекста вызывает
\texttt{@PreDestroy}. Запомните порядок двух средних шагов: сначала
объект создаётся и только потом получает зависимости, поэтому в конструкторе
нельзя рассчитывать на поля, внедряемые сеттером.

Сколько экземпляров бина создаст контейнер и как долго они проживут, задаёт
\term{область видимости} (\eng{scope}); варианты перечислены в
табл.~\ref{tab:07-2}. По умолчанию действует \texttt{singleton}, и это стоит
помнить: единственный экземпляр сервиса обслуживает все запросы одновременно,
поэтому изменяемых полей в нём быть не должно.

\begin{table}[htbp]
\caption{Области видимости бинов}
\label{tab:07-2}
\centering
\begin{tabular}{L{4.0cm}L{10.8cm}}
\toprule
Область & Сколько экземпляров \\
\midrule
\texttt{singleton} & один на весь контейнер (по умолчанию) \\
\texttt{prototype} & новый на каждое обращение к контейнеру \\
\texttt{request} & один на HTTP-запрос (только веб-приложения) \\
\texttt{session} & один на HTTP-сессию (только веб-приложения) \\
\texttt{application} & один на контекст сервлетов (только веб) \\
\bottomrule
\end{tabular}
\end{table}

Объявить бин можно тремя путями. Первый и самый частый~--- \term{стереотипные
аннотации} на классе: контейнер находит их при сканировании пакетов и
регистрирует класс сам. Второй~--- \term{фабричный метод} с аннотацией
\texttt{@Bean} внутри класса с \texttt{@Configuration}. Оба показаны в
листинге~\ref{lst:07-5}.

\begin{listing}[H]
\begin{minted}{java}
@Component  public class MyComponent { }        // общая
@Service    public class StudentService { }     // логика
@Repository public class StudentRepository { }  // данные
@Controller public class StudentController { }  // MVC

@Configuration
public class AppConfig {
    // Контейнер сам вызовет метод и зарегистрирует бин
    @Bean
    public PasswordEncoder passwordEncoder() {
        return new BCryptPasswordEncoder();
    }
}
\end{minted}
\caption{Стереотипные аннотации и фабричный метод}
\label{lst:07-5}
\end{listing}

Все четыре стереотипа~--- специализации \texttt{@Component}. Для контейнера они
почти равнозначны, но несут смысл для читателя кода, а \texttt{@Repository}
вдобавок переводит исключения конкретной СУБД в единую иерархию
\texttt{DataAccessException}. Фабричный же метод нужен потому, что пометить
\texttt{@Component} можно только свой класс: если бин делают из класса чужой
библиотеки, метод соберёт объект «под капотом» и вручит контейнеру готовый
результат. Третий путь~--- описание бинов в XML-файле
(обычно \texttt{beans.xml}), который встречается в старом коде. Каждый бин там
объявляется элементом
\texttt{<bean>}; вложенный \texttt{<constructor-arg ref="..."/>} передаёт в
конструктор другой бин, а \texttt{<property name="..." value="..."/>}~---
значение через сеттер. Загружают такой контекст строкой
\texttt{new ClassPathXmlApplicationContext("beans.xml")}.

\begin{oshibka}
\textbf{Частая ошибка: две реализации одного интерфейса.} Если пометить
\texttt{@Service} два класса, реализующих \texttt{GreetingService}, контейнер
при внедрении не поймёт, какой из них взять, и приложение упадёт на старте с
\texttt{NoUniqueBeanDefinitionException}. Лечится это тремя способами: пометить
предпочтительную реализацию \texttt{@Primary}, указать нужную по имени через
\texttt{@Qualifier} или внедрить сразу \texttt{List<GreetingService>}~---
контейнер соберёт в список все реализации.
\end{oshibka}

\section{Аспектно-ориентированное программирование}

Логирование, замер времени, проверка прав, открытие транзакции~--- всё это
нужно во множестве методов и не относится ни к одному из них по существу.
Такую функциональность называют \term{сквозной}, и если писать её руками, она
расползётся копиями по всему проекту.
\term{Аспектно-ориентированное программирование} (\eng{Aspect-Oriented
Programming, AOP})~--- парадигма, которая выносит сквозную функциональность в
отдельные модули~--- \term{аспекты}. Понятий здесь четыре. \term{Аспект}
(\eng{aspect})~--- класс с дополнительной логикой. \term{Точка соединения}
(\eng{join point})~--- место в выполнении программы, куда можно «врезаться»; в
Spring это всегда вызов метода. \term{Срез} (\eng{pointcut})~--- выражение,
отбирающее нужные точки соединения. \term{Совет} (\eng{advice})~--- код,
который выполняется в отобранной точке. Аналогия~--- турникет в метро: он
ничего не знает о том, куда вы едете, но срабатывает на каждом входе.

Аспект из листинга~\ref{lst:07-6} печатает имя каждого вызванного метода из
пакета \texttt{service}. Строка в скобках после \texttt{@Before}~--- это и есть
срез: звёздочки означают «любой тип результата», «любой класс пакета» и «любой
метод», а \texttt{(..)}~--- «с любыми параметрами». Парная аннотация
\texttt{@AfterReturning} даёт доступ к тому, что метод вернул.

\begin{listing}[H]
\begin{minted}{java}
@Aspect
@Component
public class LoggingAspect {

    @Before("execution(* com.example.service.*.*(..))")
    public void logBefore(JoinPoint jp) {
        System.out.println("Вызов: "
                + jp.getSignature().getName());
    }
}
\end{minted}
\caption{Аспект, логирующий вызовы сервисов}
\label{lst:07-6}
\end{listing}

Работает это через \term{прокси}: контейнер подменяет ваш бин
объектом-двойником с тем же интерфейсом, который сначала выполняет совет, а
потом передаёт вызов настоящему объекту. На том же механизме построены
\texttt{@Transactional}, \texttt{@PreAuthorize} и \texttt{@Cacheable}~--- то
есть, изучив AOP, вы понимаете половину «магии» Spring. Полноценная
альтернатива~--- AspectJ: эта библиотека правит байт-код при компиляции и не
ограничена вызовами методов.

\section{Spring MVC и слои приложения}

\term{Spring MVC}~--- реализация паттерна \term{Модель~--- Представление~---
Контроллер} (\eng{Model~--- View~--- Controller}). \emph{Модель}~---
бизнес-данные и логика; \emph{представление}~--- то, что видит пользователь:
HTML-страница или ответ в формате JSON; \emph{контроллер} принимает
HTTP-запрос, вызывает нужные сервисы и решает, что вернуть. На практике
приложение делят на три слоя, и запрос проходит их насквозь
(рис.~\ref{fig:07-1}). Правило простое: контроллер не лезет в базу, а
репозиторий ничего не знает про HTTP. Каждый слой заменяем независимо, и
именно поэтому веб-интерфейс приделывается к готовой бизнес-логике за час.

\begin{figure}[htbp]\centering
\begin{tikzpicture}[font=\small,node distance=4mm]
\node[boxw=16mm] (a){HTTP-\\запрос};
\node[boxw=33mm,right=of a] (b){\texttt{@RestController}\\[1pt]
  \scriptsize приём и разбор};
\node[boxw=22mm,right=of b] (c){\texttt{@Service}\\[1pt]
  \scriptsize бизнес-логика};
\node[boxw=27mm,right=of c] (d){\texttt{@Repository}\\[1pt]
  \scriptsize доступ к данным};
\node[boxw=16mm,right=of d] (e){база\\данных};
\draw[flowarrow] (a)--(b);\draw[flowarrow] (b)--(c);
\draw[flowarrow] (c)--(d);\draw[flowarrow] (d)--(e);
\end{tikzpicture}
\caption{Прохождение запроса по слоям приложения}\label{fig:07-1}
\end{figure}

Сквозной пример главы~--- сущность \texttt{Student} из
листинга~\ref{lst:07-7}: её мы отдадим через REST, покажем на HTML-странице и
сохраним в базу. Аннотации JPA знакомы по предыдущей главе.

\begin{listing}[H]
\begin{minted}{java}
@Entity
@Table(name = "students")
public class Student {

    @Id
    @GeneratedValue(strategy = GenerationType.IDENTITY)
    private Long id;

    @Column(nullable = false, length = 100)
    private String name;
    @Column(nullable = false, length = 100)
    private String surname;

    public Student() {}          // нужен Hibernate

    public Student(String name, String surname) {
        this.name = name;
        this.surname = surname;
    }
    // ... геттеры и сеттеры для всех полей
}
\end{minted}
\caption{Сущность Student}
\label{lst:07-7}
\end{listing}

Контроллеры бывают двух видов. Класс с \texttt{@Controller} возвращает
\emph{представление}: строка, которую вернул метод, воспринимается механизмом
\texttt{ViewResolver} как имя шаблона~--- \texttt{"students"} означает файл
\texttt{templates/students.html}. Класс с \texttt{@RestController}~--- это
\texttt{@Controller} плюс \texttt{@ResponseBody} на всех методах: возвращённый
объект сериализуется в JSON и уходит прямо в тело ответа. Первый вид нужен для
страниц, второй~--- для REST API.

Какой метод вызвать, контейнер решает по аннотациям сопоставления: общей
\texttt{@RequestMapping} или её сокращениям \texttt{@GetMapping},
\texttt{@PostMapping}, \texttt{@PutMapping}, \texttt{@DeleteMapping}; в
шаблонах путей звёздочка заменяет один сегмент, а две~--- любую вложенность
(\texttt{/admin/**}). Данные из запроса попадают в параметры метода четырьмя
аннотациями: \texttt{@PathVariable} берёт значение из пути,
\texttt{@RequestParam}~--- из строки запроса, \texttt{@RequestBody}
десериализует тело запроса (обычно JSON) в Java-объект, а
\texttt{@ModelAttribute} собирает объект из полей отправленной
HTML-формы. Первые три собраны в контроллере из
листинга~\ref{lst:07-8}.

\begin{listing}[H]
\begin{minted}{java}
@RestController
@RequestMapping("/api/students")
public class StudentRestController {

    // ... поле service и внедрение через конструктор

    @GetMapping
    public List<Student> getAll() {
        return service.findAll();
    }

    @GetMapping("/{id}")
    public ResponseEntity<Student> getById(@PathVariable Long id) {
        Student s = service.findById(id);
        return s != null ? ResponseEntity.ok(s)
                         : ResponseEntity.notFound().build();
    }

    @PostMapping
    public ResponseEntity<Student> create(@RequestBody Student s) {
        return ResponseEntity.status(HttpStatus.CREATED)
                             .body(service.save(s));
    }
}
\end{minted}
\caption{REST-контроллер с операциями над сущностью}
\label{lst:07-8}
\end{listing}

Класс \texttt{ResponseEntity}~--- обёртка над HTTP-ответом: он задаёт не
только тело, но и код состояния с заголовками. Соблюдайте принятые коды:
создание~--- 201~Created, удаление без тела ответа~--- 204~No~Content,
отсутствие ресурса~--- 404~Not~Found. Сам по себе Spring не возвращает 404,
когда метод отдал \texttt{null},~--- об этом заботится разработчик.

Ниже контроллера лежит сервис (класс с \texttt{@Service} и бизнес-логикой), а
ниже~--- репозиторий, и здесь Spring Data JPA избавляет от почти всей рутины:
реализацию писать не нужно, достаточно объявить интерфейс
(листинг~\ref{lst:07-9}). \texttt{CrudRepository<T, ID>} даёт методы
\texttt{save}, \texttt{findById}, \texttt{findAll}, \texttt{deleteById} и
\texttt{count}; \texttt{JpaRepository<T, ID>} добавляет сортировку,
постраничный вывод и пакетные операции.

\begin{listing}[H]
\begin{minted}{java}
public interface StudentRepository
        extends JpaRepository<Student, Long> {

    // Реализацию Spring Data напишет сама по имени метода
    List<Student> findByNameContainingIgnoreCase(String part);
}
\end{minted}
\caption{Репозиторий Spring Data JPA}
\label{lst:07-9}
\end{listing}

Метод \texttt{findByNameContainingIgnoreCase} нигде не реализован, и всё же он
работает: Spring Data разбирает имя метода по словам. Приставка
\texttt{findBy} означает выборку, \texttt{Name}~--- поле сущности,
\texttt{Containing}~--- поиск подстроки, \texttt{IgnoreCase}~--- без учёта
регистра. У правила есть обратная сторона: опечатка в имени поля обнаружится
не при компиляции, а при запуске приложения.

\section{Шаблонизатор Thymeleaf}

\term{Шаблонизатор}~--- программа, которая собирает документ из шаблона и
данных, отделяя структуру страницы от содержимого. \term{Thymeleaf}~---
современный серверный шаблонизатор для Java; его главная черта~---
\term{естественные шаблоны} (\eng{natural templates}): файл остаётся
корректным HTML, его можно открыть в браузере и увидеть вёрстку с
подставленными для примера значениями. Шаблоны кладут
в каталог \texttt{src/main/resources/templates}, статические файлы~--- в
\texttt{src/main/resources/static}. Данные контроллер передаёт через объект
\texttt{Model}: вызов \texttt{model.addAttribute("students", list)} кладёт
список под именем \texttt{students}, а возвращённое имя \texttt{"students"}
указывает на файл \texttt{students.html}. Все атрибуты Thymeleaf начинаются с
префикса \texttt{th:}~--- они читают значение из модели и подставляют его в
обычный HTML-тег (листинг~\ref{lst:07-10}). Обратите внимание на текст внутри
тегов~--- «Заголовок», «Имя»: он виден дизайнеру при открытии файла и
замещается данными при работе приложения.

\begin{listing}[H]
\begin{minted}{html}
<html xmlns:th="http://www.thymeleaf.org">
<body>

<h1 th:text="${title}">Заголовок</h1>

<table>
    <tr th:each="s : ${students}">
        <td th:text="${s.name}">Имя</td>
        <td><a th:href="@{/students/delete/{id}(id=${s.id})}">
            Удалить</a></td>
    </tr>
</table>

<form th:action="@{/students/save}" th:object="${student}"
      method="post">
    <input type="text" th:field="*{name}"/>
    <button type="submit">Сохранить</button>
</form>

</body>
</html>
\end{minted}
\caption{Шаблон Thymeleaf: вывод, цикл, ссылка и форма}
\label{lst:07-10}
\end{listing}

Атрибут \texttt{th:text} заменяет содержимое тега значением выражения,
\texttt{th:each} повторяет тег для каждого элемента коллекции, а
\texttt{th:href} с записью \texttt{@\{...\}} строит адрес. Форму связывает с
объектом пара атрибутов: \texttt{th:object} называет объект модели, а
\texttt{th:field} со звёздочкой~--- его поле; Spring при отправке формы соберёт
из них объект обратно.

\begin{vrezka}
\textbf{Почему после сохранения нужен redirect.} Метод, обрабатывающий
отправку формы, возвращает не имя шаблона, а строку
\texttt{"redirect:/students"}: браузер переходит по адресу заново, и
обновление страницы клавишей F5 больше не отправляет форму повторно~---
иначе один и тот же студент добавлялся бы при каждом нажатии.
\end{vrezka}

\section{Spring Security}

\term{Spring Security}~--- фреймворк, отвечающий за \term{аутентификацию}
(\eng{authentication}, «кто вы такой») и \term{авторизацию}
(\eng{authorization}, «что вам разрешено»), а заодно закрывающий типовые
уязвимости веб-приложений. Работает он через \term{цепочку фильтров}
(\texttt{SecurityFilterChain}): каждый входящий запрос по очереди проходит
фильтры, которые определяют пользователя, проверяют его права и отсекают
подделку запросов. Настраивают её, объявив бин
\texttt{SecurityFilterChain} (листинг~\ref{lst:07-11}). Читается конфигурация
сверху вниз, и первое подошедшее правило выигрывает~--- поэтому частные пути
пишут раньше общих, а \texttt{anyRequest} всегда идёт последним.

\begin{listing}[H]
\begin{minted}{java}
@Configuration
@EnableMethodSecurity
public class SecurityConfig {

    @Bean
    public SecurityFilterChain filterChain(HttpSecurity http)
            throws Exception {
        http
            .authorizeHttpRequests(auth -> auth
                .requestMatchers("/login").permitAll()
                .requestMatchers("/api/students")
                    .hasAnyRole("USER", "ADMIN")
                .requestMatchers("/api/students/**").hasRole("ADMIN")
                .anyRequest().authenticated()
            )
            .formLogin(form -> form.permitAll())
            .httpBasic(Customizer.withDefaults())
            .csrf(AbstractHttpConfigurer::disable);

        return http.build();
    }
}
\end{minted}
\caption{Конфигурация Spring Security}
\label{lst:07-11}
\end{listing}

Бин \texttt{PasswordEncoder} из листинга~\ref{lst:07-5} обязателен и здесь:
пароли в базе хранят только в виде хеша. Алгоритм BCrypt добавляет к паролю случайную «соль» и намеренно считает
хеш медленно, из-за чего перебор становится невыгодным, а восстановить
исходный пароль из хеша нельзя вовсе. Откуда брать пользователей, Spring
Security узнаёт через интерфейс \texttt{UserDetailsService} с методом
\texttt{loadUserByUsername}: в учебном проекте его роль играет
\texttt{InMemoryUserDetailsManager}, в настоящем~--- ваша реализация, читающая
пользователя из базы. Учебные учётные записи объявляет ещё один бин того же
класса \texttt{SecurityConfig} (листинг~\ref{lst:07-12}); пароль и здесь
хранится не текстом, а результатом \texttt{encoder.encode}.

\begin{listing}[H]
\begin{minted}{java}
@Bean
public UserDetailsService users(PasswordEncoder encoder) {
    UserDetails user = User.withUsername("user")
        .password(encoder.encode("password"))
        .roles("USER").build();
    UserDetails admin = User.withUsername("admin")
        .password(encoder.encode("password"))
        .roles("ADMIN").build();
    return new InMemoryUserDetailsManager(user, admin);
}
\end{minted}
\caption{Два пользователя в памяти приложения}
\label{lst:07-12}
\end{listing}

Два способа входа в конфигурации стоит различать. Вызов
\texttt{formLogin} рисует HTML-форму и рассчитан на браузер, а
\texttt{httpBasic}~--- на программы: клиент передаёт имя и пароль прямо в
заголовке запроса. Без \texttt{httpBasic} тот же \texttt{curl} вместо
данных получит переадресацию на форму входа, поэтому для REST API эту
строку не пропускают. Ограничивать доступ можно и по методам: аннотация
\texttt{@PreAuthorize("hasRole('ADMIN')")} проверяет выражение до вызова.
Правила цепочки и правила на методах дополняют друг друга, поэтому итог
определяет более строгая из них.

Для REST API вход по форме и хранение сессии неудобны: у клиента может не
быть браузера, а сервер приходится заставлять помнить всех вошедших.
Стандартное решение~--- \term{JWT} (\eng{JSON Web Token}): компактный
самодостаточный токен с именем пользователя, сроком действия и подписью
сервера. Работает он так: клиент отправляет логин и пароль на
\texttt{POST /api/auth/login}; сервер возвращает подписанный секретным ключом
токен; клиент добавляет к каждому запросу заголовок \texttt{Authorization} со
значением \texttt{Bearer} и самим токеном; отдельный фильтр в цепочке
проверяет подпись и срок годности. Обратите внимание: сервер не хранит
выданный токен~--- он лишь проверяет подпись, и подделать токен, не зная
ключа, невозможно.

\begin{oshibka}
\textbf{Секретный ключ в исходном коде.} Ключ подписи нельзя записывать
строковой константой в классе: любой, кто получит доступ к репозиторию, сможет
выпускать токены от имени администратора. Ключ хранят в переменной окружения
или в защищённом хранилище и подставляют в приложение при запуске.
\end{oshibka}

\section{Готовность к эксплуатации}

Всё, что было выше,~--- фундамент. Ниже~--- четыре техники, без которых
сервис не пускают в эксплуатацию.

\subsection{Проверка входных данных}

Контроллер~--- граница доверия: всё, что пришло снаружи, нужно проверить, иначе
в базу попадут пустые строки, отрицательные возрасты и несуществующие адреса
почты. Стандарт \term{Bean Validation} описывает проверки аннотациями
(табл.~\ref{tab:07-3}); реализует их Hibernate Validator.

\begin{table}[htbp]
\caption{Основные аннотации Bean Validation}
\label{tab:07-3}
\centering
\begin{tabular}{L{3.2cm}L{4.2cm}L{7.2cm}}
\toprule
Аннотация & Применяется к & Что проверяет \\
\midrule
\texttt{@NotNull} & любому типу & значение не \texttt{null} \\
\texttt{@NotBlank} & строке & не \texttt{null}, не пустая, не пробелы \\
\texttt{@Size} & строке, коллекции & длина в границах от \texttt{min} до
  \texttt{max} \\
\texttt{@Min}, \texttt{@Max} & числам & значение в границах \\
\texttt{@Email} & строке & корректный адрес почты \\
\texttt{@Valid} & вложенному объекту & запускает проверку внутри него \\
\bottomrule
\end{tabular}
\end{table}

Запомните главное: сами по себе аннотации на классе ничего не делают. Проверку
запускает \texttt{@Valid} на параметре метода контроллера, и происходит это
\emph{до} входа в метод~--- при неудаче Spring бросает
\texttt{MethodArgumentNotValidException}, а тело метода даже не начинает
выполняться.

\subsection{Разделение сущностей и объектов передачи}

В учебных примерах удобно возвращать из контроллера сущность \texttt{@Entity}
напрямую. В настоящем проекте это порождает четыре проблемы: \emph{утечку
внутреннего состояния}~--- поля \texttt{password} или \texttt{deleted} уходят
клиенту; \texttt{LazyInitializationException}~--- сериализация разворачивает
ленивую связь вне транзакции, и приложение падает; \emph{сцепление API и
базы}~--- переименовали поле в сущности, сломали всех клиентов;
\emph{циклические ссылки}~--- связь «студент~--- группа~--- список студентов»
разворачивается в бесконечный JSON. Решение~--- \term{объект передачи данных}
(\eng{Data Transfer Object, DTO}): отдельные классы для границ API, свои для
входящих данных и для исходящих. Записывают их через \texttt{record}
(листинг~\ref{lst:07-13}).

\begin{listing}[H]
\begin{minted}{java}
// Приходит от клиента и проверяется
public record StudentRequest(
    @NotBlank(message = "Имя обязательно")
    @Size(min = 2, max = 100)
    String name,

    @NotBlank String surname
) {}

// Уходит клиенту: только разрешённые поля
public record StudentResponse(Long id, String name,
                              String surname) {}

@Component
public class StudentMapper {

    // Обратный метод toEntity(StudentRequest) устроен так же
    public StudentResponse toResponse(Student s) {
        return new StudentResponse(
            s.getId(), s.getName(), s.getSurname());
    }
}
\end{minted}
\caption{Объекты передачи данных и преобразователь}
\label{lst:07-13}
\end{listing}

Контроллер после этого работает только с объектами передачи: принимает
\texttt{@Valid @RequestBody StudentRequest}, отдаёт \texttt{StudentResponse}, а
превращением занимается \texttt{StudentMapper}. Писать его руками не
обязательно: библиотека MapStruct сгенерирует преобразователь по интерфейсу
при компиляции.

\subsection{Единая обработка ошибок}

Когда проверка не прошла, клиент по умолчанию получает ответ 400 без единого
намёка на то, какое поле сломалось. Исправляет это глобальный перехватчик:
класс с аннотацией \texttt{@RestControllerAdvice} (то есть
\texttt{@ControllerAdvice} плюс \texttt{@ResponseBody}) собирает у себя
обработчики исключений для всех контроллеров сразу. Пример~---
листинг~\ref{lst:07-14}; заметьте, что формат ответа один на все случаи, а
поле \texttt{fieldErrors} заполняется только для ошибок проверки.

\begin{listing}[H]
\begin{minted}{java}
@RestControllerAdvice
public class GlobalExceptionHandler {

    public record ErrorResponse(Instant timestamp, int status,
            String error, String message,
            Map<String, String> fieldErrors) {}

    @ExceptionHandler(MethodArgumentNotValidException.class)
    public ResponseEntity<ErrorResponse> handleValidation(
            MethodArgumentNotValidException ex) {

        Map<String, String> errors = new HashMap<>();
        ex.getBindingResult().getFieldErrors().forEach(
            e -> errors.put(e.getField(),
                            e.getDefaultMessage()));

        return ResponseEntity.badRequest().body(new ErrorResponse(
            Instant.now(), 400, "Bad Request", "Validation failed",
            errors));
    }
}
\end{minted}
\caption{Глобальный обработчик исключений}
\label{lst:07-14}
\end{listing}

Точно так же добавляют остальные обработчики: метод, ловящий исключение
\texttt{EntityNotFoundException}, вернёт код 404, а общий метод для
\texttt{Exception}~--- всё остальное; Spring выбирает обработчик по наиболее
точному совпадению типа. Выигрыш очевиден: контроллеры остаются чистыми, без
блоков \texttt{try/catch}, а формат ошибок для всего API одинаков.

\subsection{Транзакции}

Аннотация \texttt{@Transactional} оборачивает вызов метода в транзакцию базы
данных: метод завершился нормально~--- изменения фиксируются, вылетело
непроверяемое исключение~--- откатываются целиком. Ставят её на
\emph{сервисном слое}. Не на контроллере: там транзакция осталась бы открытой
на время формирования JSON, удлиняя блокировки в базе. Не на репозитории: его
методы и так транзакционны по отдельности, а бизнес-операция обычно состоит из
нескольких вызовов, которые нужно объединить. Из параметров чаще всего нужен
флаг \texttt{readOnly = true}: он подсказывает Hibernate, что менять ничего не
будут, и тот отключает проверку изменений; есть ещё \texttt{propagation}
(поведение при вложенных вызовах), \texttt{isolation} (уровень изоляции) и
\texttt{rollbackFor}. Листинг~\ref{lst:07-15} показывает и типовое применение,
и классическую ловушку.

\begin{listing}[H]
\begin{minted}{java}
@Service
public class StudentService {

    // ... поле repository и внедрение через конструктор

    @Transactional(readOnly = true)
    public List<Student> findAll() {
        return repository.findAll();
    }

    // Оба сохранения либо пройдут, либо будут откачены
    @Transactional
    public void saveBoth(Student good, Student bad) {
        repository.save(good);
        if (bad.getName().isBlank()) {
            throw new IllegalStateException("Пустое имя");
        }
        repository.save(bad);
    }

    // Ловушка: вызов через this минует прокси,
    // транзакция НЕ откроется и отката не будет
    public void outer(Student good, Student bad) {
        this.saveBoth(good, bad);
    }
}
\end{minted}
\caption{Транзакции на сервисном слое и ловушка самовызова}
\label{lst:07-15}
\end{listing}

Ловушек у прокси ровно три, и все следуют из механизма, разобранного в разделе
про AOP: вызовы перехватывает объект-двойник. \emph{Первая}:
\texttt{@Transactional} на приватном методе не работает~--- двойник подменяет
только публичные методы. \emph{Вторая}: \term{самовызов}
(\eng{self-invocation}), показанный в листинге,~--- обращение через
\texttt{this} идёт напрямую в объект, минуя двойника; лечится выносом метода в
отдельный бин. \emph{Третья}: по умолчанию откат происходит только на
непроверяемых исключениях, наследниках \texttt{RuntimeException} и
\texttt{Error}; если метод бросит \texttt{IOException}, транзакция спокойно
зафиксируется~--- нужен \texttt{@Transactional(rollbackFor = Exception.class)}.

\praktikum{}

Понадобится JDK~21 или новее, Maven (подойдёт обёртка
\texttt{mvnw} из сгенерированного проекта) и любая среда разработки. Задания
выполняются подряд: каждое опирается на код предыдущего.

\subsection*{Задание 1. Проект Spring Boot и внедрение зависимостей}

Создайте на сайте \texttt{start.spring.io} проект: сборка Maven, язык Java,
Spring Boot версии~3.5, группа \texttt{mpt.it}, артефакт \texttt{student},
пакет
\texttt{lecture.seven.student}. Зависимости: Spring Web, Spring Data JPA,
Spring Security, Thymeleaf, H2 Database. Без сети тот же проект набирают
руками по листингам~\ref{lst:07-3} и~\ref{lst:07-2}. Проверьте, что
\texttt{mvnw clean compile} собирает проект, а \texttt{mvnw spring-boot:run}
запускает его: в консоли появится сообщение о старте Tomcat на порту~8080. Ответьте письменно, что находится в секции \texttt{parent} файла
\texttt{pom.xml} и почему у зависимостей не указаны версии.

Создайте пакет \texttt{greet}: интерфейс \texttt{GreetingService} с методом
\texttt{String greet(String name)}, реализацию
\texttt{EnglishGreetingService} (\texttt{@Service}) и контроллер
\texttt{GreetingController} (\texttt{@RestController}) с путём
\texttt{/api/greet/\{name\}} и конструкторным внедрением сервиса. Проверьте
результат в браузере, затем проведите три эксперимента и объясните каждый
письменно: уберите \texttt{@Service}; верните её и добавьте вторую
реализацию, тоже с \texttt{@Service}; перепишите контроллер на внедрение
через поле и объясните, почему \texttt{final} возможен только при внедрении
через конструктор.

\subsection*{Задание 2. Слои приложения и REST API}

Создайте сущность \texttt{Student} в пакете \texttt{model} по
листингу~\ref{lst:07-7}. Пропишите настройки базы
из листинга~\ref{lst:07-4}, добавьте репозиторий по образцу
листинга~\ref{lst:07-9}, интерфейс \texttt{StudentService} с методами
\texttt{findAll}, \texttt{findById}, \texttt{save}, \texttt{deleteById} и его
реализацию с аннотацией \texttt{@Service}, передающую вызовы репозиторию, а
затем контроллер по образцу листинга~\ref{lst:07-8}, дополнив его методами
обновления и удаления.

Стартер безопасности уже подключён, и Spring Boot пока закрывает все пути
сам: пользователь~--- \texttt{user}, пароль печатается в консоли при старте и
передаётся ключом \texttt{-u}. Проверьте все пять операций утилитой
\texttt{curl}, записывая тело ответа и код состояния;
создание выглядит так: \texttt{curl -u user:пароль -X POST
localhost:8080/api/students -H "Content-Type: application/json" -d
'\{"name":"Ali","surname":"Hasan"\}'}. Отдельно
запросите несуществующий идентификатор и объясните, почему приходит 404.
Добавьте метод поиска с параметром \texttt{@RequestParam String q},
вызывающий \texttt{findByNameContainingIgnoreCase}, и объясните, откуда
Spring Data берёт его реализацию. Что на самом деле попало в таблицу, видно
в веб-консоли H2 по пути \texttt{/h2-console}.

\subsection*{Задание 3. Веб-интерфейс и разграничение доступа}

Добавьте класс \texttt{StudentWebController} (\texttt{@Controller}, путь
\texttt{/students}) с тремя методами.
\texttt{String list(Model model)} кладёт в модель список студентов и пустой
объект \texttt{Student} для формы и возвращает имя шаблона;
\texttt{String save(@ModelAttribute Student s)} по пути
\texttt{/students/save} сохраняет объект и возвращает
\texttt{redirect:/students}; \texttt{String delete(@PathVariable Long id)} по
пути \texttt{/students/delete/\{id\}} делает то же самое после удаления. Создайте в каталоге \texttt{templates}
шаблон \texttt{students.html} по образцу листинга~\ref{lst:07-10}: форма
добавления, таблица со всеми студентами и ссылка удаления в каждой строке.
Откройте страницу, добавьте трёх студентов, одного удалите и объясните
письменно, что делает \texttt{th:object}, как \texttt{th:field} связывает поле
формы с полем объекта и зачем нужна приставка \texttt{redirect:}.

Добавьте в пакет \texttt{config} класс \texttt{SecurityConfig} по образцу
листингов~\ref{lst:07-11} и~\ref{lst:07-12}: цепочка фильтров, кодировщик
паролей и два пользователя в памяти с паролем \texttt{password}. Проверьте
три случая, записав коды ответов: список без входа
(\texttt{curl localhost:8080/api/students}); удаление первого студента от
имени \texttt{user} (\texttt{curl -u user:password -X DELETE
localhost:8080/api/students/1}); то же от имени \texttt{admin}. Объясните,
почему два последних кода разные: путь с идентификатором подпадает уже под
правило \texttt{/api/students/**}. Затем добавьте над
методом удаления аннотацию \texttt{@PreAuthorize("hasRole('ADMIN')")},
повторите проверку и объясните, чем ограничение через
\texttt{requestMatchers} отличается от ограничения на методе.

\subsection*{Задание 4. Доведение сервиса до рабочего состояния}

Подключите стартер \texttt{starter-validation}. Создайте пакет \texttt{dto} с
записями \texttt{StudentRequest} (проверки \texttt{@NotBlank} и
\texttt{@Size} на обоих полях) и \texttt{StudentResponse} и преобразователем
\texttt{StudentMapper} по образцу листинга~\ref{lst:07-13}.
Перепишите REST-контроллер так, чтобы он принимал
\texttt{@Valid @RequestBody StudentRequest} и возвращал только
\texttt{StudentResponse}; отправьте с ключом \texttt{-u user:password} запрос
с пустым именем и убедитесь, что приходит код 400. Добавьте пакет
\texttt{exception} с классом \texttt{GlobalExceptionHandler} по образцу
листинга~\ref{lst:07-14} и обработчиком для
\texttt{jakarta.persistence.EntityNotFoundException}, а метод
\texttt{findById} исправьте так, чтобы при отсутствии студента он бросал это
исключение. Повторите оба запроса: теперь в первом случае придёт
структурированный JSON с полем \texttt{fieldErrors}, во втором~--- код 404 с
понятным сообщением.

Расставьте в сервисе аннотации \texttt{@Transactional}: с флагом
\texttt{readOnly = true} на методах чтения и без него на изменяющих. Добавьте
метод \texttt{saveBoth} из листинга~\ref{lst:07-15}. Снаружи его вызывают
временным эндпоинтом: метод контроллера с \texttt{@GetMapping("/demo-tx")},
внутри которого стоит
\texttt{saveBoth(new Student("Ali", "Hasan"), new Student("", "X"))}.
Откройте этот адрес и проверьте по списку, что первый студент \emph{не}
сохранился; уберите аннотацию, повторите вызов и сравните результат. Затем
воспроизведите ловушку самовызова из того же листинга и объясните, почему
теперь первый студент сохраняется. По желанию подключите стартер
\texttt{starter-aop} и добавьте аспект логирования по образцу
листинга~\ref{lst:07-6}. Ответьте письменно, что произойдёт, если убрать
\texttt{@Valid} с параметра.

\subsection*{Результаты занятия}

По итогам занятия должен быть готов проект Spring Boot: REST-контроллер с
полным набором операций над студентами, веб-интерфейс на Thymeleaf,
безопасность с двумя учётными записями и разграничением прав, объекты
передачи данных с проверками и преобразователь, глобальный обработчик ошибок,
транзакции на сервисном слое. Дополнительно~--- письменные ответы на вопросы
всех четырёх заданий и записанные коды ответов.

\kontrolvoprosy

\begin{enumerate}
\item Сформулируйте принцип инверсии управления своими словами. Чем он
      отличается от подхода, в котором объект сам создаёт свои зависимости?
\item Назовите три способа внедрения зависимостей. Какой предпочтителен и по
      каким причинам?
\item Что даёт аннотация \texttt{@SpringBootApplication}, почему главный класс
      кладут в корневой пакет и что такое starter-модуль?
\item Перечислите шаги жизненного цикла бина. Почему в конструкторе нельзя
      рассчитывать на зависимость, внедрённую через сеттер?
\item Какие области видимости бинов вы знаете, чем опасно изменяемое поле в
      бине с областью \texttt{singleton} и когда бин нельзя объявить
      аннотацией \texttt{@Component}?
\item Объясните понятия «аспект», «точка соединения», «срез» и «совет». Через
      какой механизм AOP работает в Spring?
\item Чем \texttt{@RestController} отличается от \texttt{@Controller} и чем
      \texttt{@PathVariable} отличается от \texttt{@RequestParam}?
\item Зачем приложение делят на слои «контроллер~--- сервис~--- репозиторий»?
      Как Spring Data строит запрос по имени метода вида
      \texttt{findByNameContainingIgnoreCase}?
\item Что такое \texttt{SecurityFilterChain} и в каком порядке просматриваются
      правила доступа? Опишите порядок работы аутентификации по токену JWT.
\item Назовите четыре проблемы, возникающие при возврате сущности прямо из
      REST-контроллера, и объясните, почему самовызов
      \texttt{@Transactional}-метода не открывает транзакцию.
\end{enumerate}

\testzadaniya

\begin{enumerate}

\vopros{Что такое инверсия управления?}
  {способ оптимизации SQL-запросов}
  {принцип, при котором управление жизненным циклом объектов передаётся
   внешнему контейнеру}
  {шаблон проектирования для организации многопоточности}
  {изменение порядка вызовов методов в стеке}

\vopros{Почему внедрение через поле считают плохой практикой?}
  {оно работает медленнее конструкторного}
  {Spring его больше не поддерживает}
  {поле нельзя сделать \texttt{final}, зависимости скрыты, тест без
   контекста требует рефлексии}
  {оно требует обязательной XML-конфигурации}

\vopros{Что произойдёт, если две реализации одного интерфейса помечены
  \texttt{@Service} и ничем не различены?}
  {Spring возьмёт первую попавшуюся}
  {Spring внедрит список из обеих}
  {приложение упадёт с \texttt{NoUniqueBeanDefinitionException}}
  {приложение запустится, но бин не будет создан}

\vopros{Что включает в себя аннотация \texttt{@SpringBootApplication}?}
  {\texttt{@RestController}, \texttt{@Entity}, \texttt{@Repository}}
  {\texttt{@Configuration}, \texttt{@EnableAutoConfiguration},
   \texttt{@ComponentScan}}
  {\texttt{@Service}, \texttt{@Component}, \texttt{@Bean}}
  {\texttt{@SpringBootTest} и \texttt{@AutoConfigureMockMvc}}

\vopros{Что описывает понятие «срез» в AOP?}
  {класс с дополнительной логикой}
  {код, выполняемый при срабатывании аспекта}
  {точку остановки в отладчике}
  {выражение, отбирающее точки соединения для перехвата}

\vopros{Чем \texttt{@RestController} отличается от \texttt{@Controller}?}
  {он работает только с XML}
  {он предназначен для статических ресурсов}
  {это \texttt{@Controller} плюс \texttt{@ResponseBody}: результат
   сериализуется в тело ответа}
  {он не поддерживает запросы GET}

\vopros{Как клиент передаёт токен JWT при обращении к защищённому ресурсу?}
  {в параметре строки запроса}
  {в заголовке \texttt{Authorization} со словом \texttt{Bearer}}
  {в файле cookie}
  {в теле каждого запроса отдельным полем}

\vopros{Чем \texttt{@NotBlank} отличается от \texttt{@NotNull} для строки?}
  {они полностью равнозначны}
  {\texttt{@NotBlank} запрещает \texttt{null}, пустую строку и строку из
   одних пробелов}
  {\texttt{@NotBlank} применяется только к коллекциям}
  {\texttt{@NotBlank} проверяет совпадение с регулярным выражением}

\vopros{Где правильно ставить \texttt{@Transactional}?}
  {на контроллере, чтобы транзакция жила весь запрос}
  {на репозитории, иначе он не работает}
  {на главном классе приложения}
  {на сервисном слое, где заключена бизнес-логика}

\vopros{При каких исключениях транзакция откатывается по умолчанию?}
  {при любых, включая проверяемые}
  {только при ошибках базы данных}
  {никогда: откат вызывают вручную}
  {только при непроверяемых~--- наследниках \texttt{RuntimeException} и
   \texttt{Error}}

\end{enumerate}
